\Article{11}
Die Linie auf der Erdoberfläche, auf der $V$ in jedem Punkt denselben bestimmten Wert $=V^0$ hat, trennt im Allgemeinen die Teile dieser Fläche, in denen $V$ größer als $V^0$ ist, von denen, in denen es kleiner ist. Die horizontale magnetische Kraft steht in jedem Punkt dieser Linie offenbar senkrecht auf ihr und ist zu der Seite hin gerichtet, auf der die größeren Werte von $V$ auftreten. Ist $\mathrm{d}s$ eine unendlich kleine Linie in dieser Richtung und $V^0+\mathrm{d}V^0$ der Wert von $V$ an ihrem anderen Endpunkt, so ist $\frac{\mathrm{d}V^0}{\mathrm{d}s}$ die Intensität der horizontalen magnetischen Kraft an dieser Stelle. Auch hier bilden alle Punkte, die dem Wert $V=V^0+\mathrm{d}V^0$ entsprechen, eine zweite Linie, die der ersten unendlich nahe liegt. Dadurch wird auf der Erdoberfläche eine \textit{Zone} abgegrenzt, in der die Werte von $V$ zwischen $V^0$ und $V^0+\mathrm{d}V^0$ liegen und die horizontale Intensität umgekehrt proportional zur Breite der Zone ist. Lässt man also $V$ in unendlich kleinen, aber gleichen Schritten vom kleinsten auf der Erdoberfläche vorkommenden Wert bis zum größten anwachsen, so wird die gesamte Erdoberfläche in unendlich viele unendlich schmale Zonen geteilt. Die horizontale magnetische Kraft steht überall senkrecht auf ihren Trennlinien, und ihre Intensität ist umgekehrt proportional zur Breite der Zonen an den jeweiligen Stellen. Den beiden äußersten Werten von $V$ entsprechen dabei zwei von den Zonen eingeschlossene Punkte, an denen die horizontale Kraft $=0$ wird und \SourceLeaf{138}die gesamte magnetische Kraft folglich nur senkrecht sein kann. Diese Punkte heißen die magnetischen Pole der Erde.

Die Trennlinien der Zonen sind die Schnittlinien der in Artikel 7 betrachteten Flächen mit der Erdoberfläche. An den Polen berühren diese Flächen die Erdoberfläche lediglich.

\Article{12}
Die im vorigen Artikel beschriebene Form des Liniensystems ist eigentlich nur der einfachste Typ. Er könnte vielerlei Ausnahmen zulassen, wenn man jede mögliche Verteilung des Magnetismus in der Erde berücksichtigte. Wir werden diesen Gegenstand hier aber nicht erschöpfend behandeln, sondern zur Erläuterung nur einige Bemerkungen über die Ausnahmefälle anfügen. zumal das Liniensystem auf der Erdoberfläche bei der \textit{wirklichen} magnetischen Beschaffenheit der Erde durchaus jene Form hat. Jedenfalls gibt es gewiss keine großräumigen Ausnahmefälle, sondern höchstens vielleicht hier und da einen örtlich begrenzten.

Einige Physiker haben die Ansicht vertreten, dass die Erde zwei magnetische Nordpole und zwei Südpole habe. Die wesentlichste Bedingung scheint aber zuvor unerfüllt geblieben zu sein. Es wurde keine \textit{genaue} Begriffsbestimmung dafür gegeben, was unter einem magnetischen Pol zu verstehen sei. Wir bezeichnen mit diesem Namen jeden Punkt der Erdoberfläche, an dem die horizontale Intensität $=0$ ist. Im Allgemeinen ist dort also die Inklination $=90^0$. Eingeschlossen ist aber auch der Sonderfall, falls er auftreten sollte, in dem die gesamte Intensität $=0$ ist. Wollte man die Stellen als magnetische Pole bezeichnen, an denen die gesamte Intensität einen Höchstwert hat, also größer ist als überall in der unmittelbaren Umgebung, so darf man nicht vergessen, dass dies von jener Begriffsbestimmung ganz verschieden ist. Diese Stellen hängen weder ihrer Lage noch ihrer Anzahl nach notwendig mit jenen Punkten zusammen, und es führt zu Verwirrung, wenn man verschiedenartige Dinge gleich benennt.

Sieht man von der wirklichen Beschaffenheit der Erde ab und fasst die Frage allgemein auf, so können durchaus mehr als zwei magnetische Pole existieren. Es scheint aber bisher unbemerkt geblieben zu sein, dass etwa bei zwei Nordpolen zwischen ihnen noch ein dritter Punkt liegen \SourceLeaf{139}muss, der ebenfalls ein magnetischer Pol ist, eigentlich weder Nordpol noch Südpol, oder, wenn man lieber will, beides zugleich.

Am besten lässt sich dies durch Betrachtung unseres Liniensystems klären.

\expandafter\def\csname HistoricalNoteText16\endcsname{Auch wenn man diese Ringlinien als unendlich klein annimmt, sind sie nicht notwendig kreisrund, sondern im Allgemeinen elliptisch. Daher fällt die zu ihnen senkrechte Richtung der Magnetnadel nur an vier Stellen jedes Rings mit der Richtung nach $P^*$ zusammen. Man kann daher beträchtliche Fehler machen, wenn man den Schnittpunkt zweier aus beträchtlicher Entfernung verlängerter Kompassrichtungen ohne Weiteres für $P^*$ hält.}
\expandafter\def\csname HistoricalNoteText17\endcsname{Wir folgen hier dem üblichen Sprachgebrauch, der den von Kapitän Ross bestimmten Punkt so nennt, obwohl er eigentlich ein Südpol ist, sofern man die Erde selbst als Magneten betrachtet.}
Hat die Funktion $V$ in einem Punkt $P^*$ der Erdoberfläche einen Höchstwert $V^*$, ringsum also kleinere Werte, so entspricht einer Reihe stufenweise abnehmender Werte ein System von Ringlinien. Jede schließt alle vorhergehenden und den Punkt $P^*$ ein. Auf jeder dieser Ringlinien weist die horizontale magnetische Kraft oder der Nordpol der Magnetnadel \textit{nach innen}\HistoricalNote{16}. Dies ist das Kennzeichen eines magnetischen Nordpols{\renewcommand{\thefootnote}{**)}\HistoricalNote{17}}. Offenbar kann man die Ringe so klein wählen oder die zugehörigen Werte der Funktion $V$ so wenig von $V^*$ abweichen lassen, dass jeder andere gegebene Punkt noch außerhalb liegt.

Wir wollen mit $S$ die Gesamtheit aller Punkte der Erdoberfläche bezeichnen, in denen der Wert von $V$ größer ist als eine gegebene Größe $W$. Offenbar bildet $S$ entweder einen einzigen zusammenhängenden Teil der Fläche oder mehrere voneinander getrennte Teile. \textit{Auf} der Begrenzungslinie oder den Begrenzungslinien, die sie von den übrigen Teilen trennen, in denen $V$ kleiner als $W$ ist, gilt $V=W$. Lässt man $W$ ab- oder zunehmen, so weitet sich dieser Flächenbereich aus oder wird enger.

Nehmen wir nun an, $P^{**}$ sei ein zweiter Punkt von gleicher Art wie $P^*$, sodass auch dort $V$ einen Höchstwert $=V^{**}$ hat. Wie bereits bemerkt, kann man der Größe $W$ einen Wert zuweisen, der kleiner als $V^*$ ist und so wenig davon abweicht, dass $P^{**}$ außerhalb desjenigen Teils von $S$ liegt, in dem $P^*$ liegt. Setzt man nun voraus, dass $V^{**}$ nicht kleiner als $V^*$ ist, was zulässig ist, und damit auch größer als $W$, so muss \SourceLeaf{140}auch $P^{**}$ zu einem Teil von $S$ gehören. Folglich liegen $P^*$ und $P^{**}$ beide in $S$, aber in getrennten Teilen von $S$.

Andererseits kann man $W$ offenbar so klein wählen, dass $P^*$ und $P^{**}$ in einem einzigen zusammenhängenden Teil von $S$ liegen. Wählt man nämlich $W$ nur klein genug, so kann $S$ die gesamte Erdoberfläche umfassen.

Lässt man nun $W$ alle Werte vom ersten zum zweiten stufenweise durchlaufen, so muss es darunter einen letzten Wert $V^{***}$ geben, für den $P^*,P^{**}$ noch in getrennten Teilen von $S$ liegen. Sobald $W$ von dort aus weiter abnimmt, fließen diese Teile zu einem zusammen.

Geschieht dieses Zusammenfließen in einem einzigen Punkt $P^{***}$. so hat die Begrenzungslinie, auf der $V=V^{***}$ gilt, die Form einer Acht, die sich in diesem Punkt selbst kreuzt. Man sieht leicht, dass dort die horizontale Intensität $=0$ sein muss. Die Kreuzung erfolgt nämlich entweder unter einem messbaren Winkel oder nicht. Im ersten Fall müsste die horizontale Kraft, falls sie nicht $=0$ wäre, senkrecht auf zwei verschiedenen Tangenten stehen, was widersinnig ist. Im zweiten Fall würden sich die beiden Hälften der Acht in $P^{***}$ berühren oder dieselbe Tangente haben. Die zu dieser Tangente senkrechte Kraft könnte dann nur ins Innere einer der beiden von der Acht umschlossenen Flächenhälften gerichtet sein. Das ist ein Widerspruch, da der Wert von $V$ zu beiden Seiten hin wächst. Damit ist $P^{***}$ nach unserer Definition ein wirklicher magnetischer Pol. In Bezug auf die nächstgelegenen Punkte innerhalb der beiden Schleifen der Acht muss er als Südpol, in Bezug auf die außerhalb liegenden Punkte dagegen als Nordpol betrachtet werden. Figur 1 kann diese Form des Liniensystems veranschaulichen.

Geschieht das Zusammenfließen an zwei verschiedenen Stellen zugleich, so gilt für sie dasselbe, was eben für einen einzelnen Punkt bewiesen wurde. Man sieht leicht, dass sich innerhalb des Bereichs, der $P^*$ und $P^{**}$ einschließt, dann ein inselförmiger Bereich bildet. Bei fortgesetzter Abnahme von $W$ wird er immer enger und muss zuletzt zu einem wirklichen Südpol zusammenschrumpfen.

Ähnliches gilt, wenn das Zusammenfließen zugleich in drei oder mehr einzelnen Punkten stattfindet. Geschieht es aber auf einmal entlang einer ganzen Linie, so muss auch in allen ihren Punkten die horizontale Kaft verschwinden.

\SourceLeaf{141}Übrigens ist unmittelbar klar, dass ebenso die Annahme von zwei \textit{Südpolen} die Existenz eines dritten Polpunkts voraussetzt, der weder Südpol noch Nordpol, oder vielmehr beides zugleich ist.

\Article{13}
Aus dem, was im vorigen Artikel entwickelt wurde, erkennt man nun leicht. wie mehrere denkbare Ausnahmen vom einfachsten Typ unseres Liniensystems zustande kommen können. Die Gesamtheit aller Punkte, denen ein bestimmter Wert von $V$ entspricht, kann eine Linie sein, die aus mehreren Teilen besteht. Jeder von ihnen kehrt in sich zurück, aber alle sind vollständig voneinander getrennt. Es kann eine Linie sein, die sich selbst kreuzt. Schließlich kann es auch eine Linie sein, an die auf beiden Seiten Flächenbereiche grenzen, in denen $V$ größer ist als auf der Linie, oder auf beiden Seiten kleiner.

Wir können behaupten, dass großräumige Abweichungen dieser Art vom einfachsten Typ auf der Erde nicht vorkommen. Örtliche Abweichungen sind aber gut denkbar, wenn sich dicht unter der Erdoberfläche magnetische Massen befinden. In einiger Entfernung üben sie keine merkliche Wirkung mehr aus, in der unmittelbaren Umgebung aber eine so starke, dass sie die regelmäßig verlaufende erdmagnetische Kraft völlig übertreffen und unkenntlich machen. Im einfachsten Fall könnte das Liniensystem in einer solchen Gegend die Form haben, die die $2^{\mathrm{te}}$ Figur veranschaulicht.

